\section{مدیریت اعتماد و شهرت}

\subsection{مفهوم اعتماد در شبکه اقتضایی خودرویی}

فصل شش نشان داد که سازوکارهای رمزنگاری می‌توانند اصالت، یکپارچگی و انکارناپذیری پیام‌ها را تضمین کنند، اما این تضمین‌ها به صحت محتوای پیام تعمیم نمی‌یابد. در الگوی تهدید معرفی‌شده در فصل ۴، دسته‌ای از مهاجمان «گره‌های موذی داخلی» هستند: گره‌هایی که گواهی معتبر دارند و پیام‌های خود را از نظر رمزنگاری صحیح امضا می‌کنند، اما محتوای پیام‌هایشان عمداً نادرست است. رفتارهای خودخواهانه، تبانی بین گره‌ها و انتشار اطلاعات نادرست نیز از همین دسته محسوب می‌شوند \cite{farsimadanReviewSecurityChallenges2025}. به عبارت دقیق‌تر، سازوکارهای رمزنگاری پیشگیرانه در برابر حملات داخلی گره‌های مجاز، مانند جعل موقعیت و تولید هشدار نادرست، آسیب‌پذیرند \cite{farsimadanReviewSecurityChallenges2025}.

در بسیاری از سازوکارهای متمرکز، واحدهای کنارجاده‌ای به‌عنوان عنوان کاملاً قابل اعتماد فرض می‌شوند، در حالی که همیشه در دسترس نیستند و خودشان آسیب‌پذیرند \cite{farsimadanReviewSecurityChallenges2025}. ازاین‌رو در شبکه‌های خودرویی که هیچ مرجع متمرکزی نمی‌تواند پیوسته رفتار گره‌ها را داوری کند، اعتماد به معنای ارزیابی قابلیت اتکای یک گره است. در این شبکه‌ها اعتماد به‌عنوان اطمینان یک عنصر نسبت به عنصر دیگری تعریف شده که مؤلفه‌ای از شبکهٔ خودرویی باشد و محاسبهٔ آن یک معیار امنیتی تکمیلی محسوب می‌شود \cite{farsimadanReviewSecurityChallenges2025}؛ مدیریت اعتماد نیز به مجموعهٔ سازوکارهایی اشاره دارد که یک خودرو با آن‌ها قابلیت اطمینان سایر خودروها و ارتباطات دریافتی را ارزیابی می‌کند \cite{farsimadanReviewSecurityChallenges2025}.

مدیریت اعتماد و شهرت به‌عنوان راهکاری مکمل رمزنگاری معرفی می‌شود که با بررسی رفتار گره‌ها و محاسبهٔ امتیاز اعتماد، ارزیابی قابلیت اتکای گره‌ها را ممکن می‌سازد \cite{razaBlockchainBasedReputationTrust2024}. از دیدگاه مرور نظام‌مند رضا و همکاران، مدیریت اعتماد از گردآوری اطلاعات، نگهداری محتوای مرتبط با اعتماد، محاسبهٔ اعتماد، نگهداری اعتماد و تصمیم‌گیری تشکیل شده و خط دفاع اول در ابعاد امنیتی سامانه‌های دیجیتال به شمار می‌آید \cite{razaBlockchainBasedReputationTrust2024}. مدل کلی آن شامل ترکیب، انتشار، تجمیع، به‌روزرسانی و شکل‌گیری است که به قابلیت اتکا می‌انجامد \cite{razaBlockchainBasedReputationTrust2024}. همچنین سامانه‌های مدیریت اعتماد نقشی اساسی در حفظ حریم خصوصی و هویت خودروها و داده‌های آن‌ها دارند \cite{razaBlockchainBasedReputationTrust2024}، و ناکارآمدی، تأخیر، ارزیابی اعتماد ناهمگن و تجمیع بر گره‌های بسیار پویا از مشهورترین مسائل مدیریت اعتماد در اینترنت خودروها\LTRfootnote{\lr{Internet of Vehicles (IoV)}} شناخته شده است \cite{razaBlockchainBasedReputationTrust2024}.

\subsection{مدل‌های مدیریت اعتماد}

پس از تعریف اعتماد، پرسش این است که چه مدل‌هایی برای مدیریت آن طرح شده است. مرور رضا و همکاران سامانه‌های موجود مدیریت اعتماد را از حیث رویکرد به چهار ردهٔ مبتنی بر توصیه\LTRfootnote{\lr{Recommendation-based}}، مبتنی بر پیش‌بینی\LTRfootnote{\lr{Prediction-based}}، مبتنی بر شهرت\LTRfootnote{\lr{Reputation-based}} و مبتنی بر سیاست\LTRfootnote{\lr{Policy-based}} دسته‌بندی می‌کند \cite{razaBlockchainBasedReputationTrust2024}؛ این مرور برای هر یک از این رده‌ها تعریف تفصیلی ارائه نمی‌دهد. از حیث روش، مدل‌های تصمیم‌گیری، ارزیابی و مدیریت شناسایی شده‌اند \cite{razaBlockchainBasedReputationTrust2024}.

از حیث حالت مدیریت، سامانه‌های مدیریت اعتماد به سه گروه متمرکز، غیرمتمرکز و نیمه‌متمرکز بخش می‌شوند \cite{razaBlockchainBasedReputationTrust2024}. در مدل متمرکز، کنترل آسان‌تر و هزینهٔ کمتر امکان‌پذیر است \cite{ghovanlooyghajarSBTMSScalableBlockchain2021}، اما این مدل گلوگاه ایجاد می‌کند، تصمیم‌گیری را کند می‌سازد و نقطهٔ شکست واحد\LTRfootnote{\lr{Single Point of Failure (SPoF)}} دارد؛ اگر سامانهٔ اصلی از کار بیفتد، کل شبکه از کار می‌افتد \cite{ghovanlooyghajarSBTMSScalableBlockchain2021}. در مدل غیرمتمرکز نقطهٔ شکست واحد حذف می‌شود، اما سربار سنگین محاسباتی و انرژی به وجود می‌آید \cite{razaBlockchainBasedReputationTrust2024}. مدل نیمه‌متمرکز هر دو را ترکیب می‌کند و بیشتر معایب هر دو را برطرف می‌سازد \cite{razaBlockchainBasedReputationTrust2024}. در این دسته‌بندی، بلاکچین و رایانش لبه\LTRfootnote{\lr{Edge/Fog Computing}} در زمرهٔ مدل‌های غیرمتمرکز و ترکیب با فناوری دفترکل توزیع‌شده\LTRfootnote{\lr{Distributed Ledger Technology (DLT)}} در زمرهٔ مدل‌های نیمه‌متمرکز قرار می‌گیرند \cite{razaBlockchainBasedReputationTrust2024}.

در مقابل، در شبکه‌های همتابه‌همتا\LTRfootnote{\lr{Peer-to-Peer (P2P)}} مدل‌های متعددی برای به‌روزرسانی باور یک گره بر اساس مقدار اعتماد گره‌های دیگر به کار رفته است؛ نمونه‌ای از این مدل‌ها، سامانهٔ \lr{TrustVote}\LTRfootnote{\lr{TrustVote}} بر پایهٔ جمع‌آوری داده از مشارکت‌کنندگان\LTRfootnote{\lr{crowdsourcing}} و رمزنگاری همومورفیک\LTRfootnote{\lr{homomorphic encryption}} است \cite{ghovanlooyghajarSBTMSScalableBlockchain2021}. در برخی طرح‌های مرورشده نیز اعتماد مستقیم و اعتماد توصیه‌ای تمایز داده شده و اعتماد سراسری به‌صورت میانگین وزنی اعتمادهای محلی تعریف شده است \cite{razaBlockchainBasedReputationTrust2024}.

\subsection{محاسبه و به‌روزرسانی اعتماد}

مدل‌های معرفی‌شده نیازمند روش‌هایی برای محاسبه و به‌روزرسانی مقدار اعتماد هستند. روش‌های محاسبهٔ اعتماد در مرور رضا و همکاران شامل استنباط بیزی، استدلال فازی، بیشینهٔ درست‌نمایی، میانگین وزنی و استدلال خاکستری است \cite{razaBlockchainBasedReputationTrust2024}؛ تکنیک‌های پیاده‌سازی سامانه‌های مدیریت اعتماد نیز شامل تولید گواهی، رمزنگاری (مانند \lr{ECC}\LTRfootnote{\lr{Elliptic Curve Cryptography (ECC)}} و \lr{ECDSA}\LTRfootnote{\lr{Elliptic Curve Digital Signature Algorithm (ECDSA)}})، استنباط بیزی، منطق فازی و مقادیر اعتماد (رتبه‌بندی یا شهرت) می‌شود \cite{razaBlockchainBasedReputationTrust2024}. رایج‌ترین راهبرد مدیریت اعتماد، جمع وزنی امتیازات اعتماد\LTRfootnote{\lr{Weighted-sum trust}} است که در آن با رفتار مخرب، اعتماد به تدریج تا صفر کاهش می‌یابد \cite{farsimadanReviewSecurityChallenges2025}.

در طرح سامانهٔ مدیریت اعتماد مقیاس‌پذیر بلاکچینی برای شبکه‌های اقتضایی خودرویی\LTRfootnote{\lr{Scalable Blockchain Trust Management System (SBTMS)}}، خودروی گیرنده با استنباط بیزی قابلیت اطمینان پیام را محاسبه می‌کند؛ این محاسبه بر دو عامل استوار است: فاصلهٔ فرستنده تا رویداد، که پیام خودروهای نزدیک‌تر به رویداد را قابل اعتمادتر می‌سازد، و سطح اعتماد فرستنده برای آن نوع پیام \cite{ghovanlooyghajarSBTMSScalableBlockchain2021}. خودرو نتیجه را به نزدیک‌ترین واحد کنارجاده‌ای می‌فرستد و آن واحد قابلیت اطمینان خالص تجمیعی را با تابع سیگموید\LTRfootnote{\lr{sigmoid function}} محاسبه می‌کند \cite{ghovanlooyghajarSBTMSScalableBlockchain2021}. قابلیت اطمینان پیام و قابلیت اطمینان خالص با گذشت زمان همگرا و افزایشی هستند؛ افزایش نرخ رویدادها این مقادیر را کاهش و افزایش شعاع ارسال و تعداد خودروها آن‌ها را افزایش می‌دهد \cite{ghovanlooyghajarSBTMSScalableBlockchain2021}. در این طرح، تراکنش زمانی تأیید می‌شود که قابلیت اطمینان خالص از آستانهٔ ۷۰ درصد بگذرد \cite{ghovanlooyghajarSBTMSScalableBlockchain2021}.

منطق فازی نیز برای محاسبهٔ اعتماد به کار رفته است؛ در یکی از طرح‌های مرورشده، تکنیک محاسبهٔ اعتماد مبتنی بر منطق فازی برای ارزیابی دقت و یکپارچگی پیام‌های رویداد و فرستنده‌های آن‌ها به‌کار گرفته شده است \cite{farsimadanReviewSecurityChallenges2025}. یادگیری تقویتی نیز در یک رویکرد مبتنی بر یادگیری \lr{Q}\LTRfootnote{\lr{Q-learning}} برای ارزیابی قابلیت اطمینان خودروهای خودران و شناسایی نفوذگران بر اساس سطح اطمینان محاسبه‌شده به کار رفته است که دو روش ارزیابی مستقیم و غیرمستقیم را در پیش دارد \cite{farsimadanReviewSecurityChallenges2025}.

گره‌های مخرب پیشرفته با جابه‌جایی میان رفتار مخرب و عادی خود را پنهان می‌سازند \cite{farsimadanReviewSecurityChallenges2025}. در طبقه‌بندی حملات مبتنی بر الگوی رفتاری، حملات خودخواهانه (مانند استراق سمع منفعل و انکار\LTRfootnote{\lr{repudiation}}) و حملات بدخواهانه (مانند حملهٔ سیاه‌چاله\LTRfootnote{\lr{blackhole attack}} و سیبل\LTRfootnote{\lr{Sybil attack}}) شناسایی می‌شوند \cite{farsimadanReviewSecurityChallenges2025}؛ در این طبقه‌بندی، دفاع از حملهٔ سیاه‌چاله با سامانهٔ تشخیص نفوذ ترکیبی و مسیریابی امن پیشنهاد شده و از منطق فازی برای تشخیص رهاکردن بسته نیز به کار رفته است \cite{farsimadanReviewSecurityChallenges2025}.

\subsection{سیستم‌های شهرت}

محاسبهٔ اعتماد در مقیاس شبکهٔ خودرویی نیازمند زیرساختی است که امتیازات و سوابق رفتاری گره‌ها را به‌صورت ایمن و تغییرناپذیر نگهداری کند؛ سامانه‌های شهرت\LTRfootnote{\lr{reputation system}} به این نیاز پاسخ می‌دهند. در برخی طرح‌ها، مقادیر اعتماد در بلاک‌ها ذخیره می‌شوند \cite{razaBlockchainBasedReputationTrust2024} و مقادیر اعتماد خودروها در بلاکچین نگهداری و میان واحدهای کنارجاده‌ای به اشتراک گذاشته می‌شوند \cite{farsimadanReviewSecurityChallenges2025}؛ به عنوان نمونه، خودروها اعتبار پیام‌های دریافتی را رتبه‌بندی کرده و به واحد کنارجاده‌ای می‌فرستند و این واحد آفست اعتماد\LTRfootnote{\lr{trust offset}} را محاسبه می‌کند \cite{razaBlockchainBasedReputationTrust2024} و اجماع ترکیبی اثبات کار\LTRfootnote{\lr{Proof of Work (PoW)}} و اثبات سهم\LTRfootnote{\lr{Proof of Stake (PoS)}} برای آن به کار رفته است \cite{razaBlockchainBasedReputationTrust2024}.

سامانهٔ \lr{DIVA}\LTRfootnote{\lr{DIVA}} شهرت را مبتنی بر شناسهٔ غیرمتمرکز\LTRfootnote{\lr{Decentralized Identifier (DID)}} بنا می‌کند: مرجع مورد اعتماد شناسهٔ غیرمتمرکز هر خودرو را بر دفتر کل ثبت می‌کند و برای آن اعتبارنامهٔ تأییدپذیر\LTRfootnote{\lr{Verifiable Credential (VC)}} صادر می‌کند، و خودرو با پیوستن ارائهٔ تأییدپذیر\LTRfootnote{\lr{Verifiable Presentation (VP)}} به هر پیام از مالکیت شناسهٔ خود پشتیبانی می‌کند \cite{feraudoDIVADIDbasedReputation2024}. امتیاز شهرت بر گره‌های لبه\LTRfootnote{\lr{edge node}} محاسبه می‌شود، نه بر خودروها، و بر بستر \lr{IOTA Tangle}، دفتر کل مبتنی بر گراف جهت‌دار بی‌دور\LTRfootnote{\lr{Directed Acyclic Graph (DAG)}} با دسترسی محدود (\lr{permissioned}\LTRfootnote{\lr{permissioned}}) که بر گره‌های لبه‌ای متصل به هستهٔ \lr{5G}\LTRfootnote{\lr{5G core}} توزیع شده، نگهداری می‌شود \cite{feraudoDIVADIDbasedReputation2024}.

الگوریتم \lr{DIVA} پیام‌های اعلان محیطی غیرمتمرکز\LTRfootnote{\lr{Decentralized Environmental Notification Message (DENM)}} را با دادهٔ واحد کنارجاده‌ای، پیام‌های آگاهی همیارانه\LTRfootnote{\lr{Cooperative Awareness Message (CAM)}} همان فرستنده و رویدادهای مشابه سنجیده و امتیاز را به‌روز می‌کند \cite{feraudoDIVADIDbasedReputation2024}. امتیاز نهایی با تراکنش بدون ارزش\LTRfootnote{\lr{zero-value transaction}} در تنگل ذخیره می‌شود؛ چنین تراکنش‌هایی به اعتبارسنجی شرکت‌کنندگان شبکه نیاز ندارند \cite{feraudoDIVADIDbasedReputation2024}. در \lr{DIVA}، خودرو تنها در صورتی می‌تواند شهرت سایر مشارکت‌کنندگان را دریافت کند که شهرت خودش از آستانه‌ای معین بگذرد \cite{feraudoDIVADIDbasedReputation2024}.

در ارزیابی شبیه‌سازی‌شدهٔ \lr{DIVA} (سناریوی تقاطع اینگولشتات نورد در آلمان با ۲۰ تا ۴۰ درصد منابع مخرب)، با آستانهٔ میانگین و ضریب \lr{β} بیش از ۰٫۳، دقت حدود ۹۴ درصد با ۳۰ درصد منبع مخرب و ۸۹ درصد با ۴۰ درصد گزارش شده است \cite{feraudoDIVADIDbasedReputation2024}؛ تأخیر انتشار شهرت نیز در مرتبهٔ چند میکروثانیه اندازه‌گیری شد \cite{feraudoDIVADIDbasedReputation2024}. باید توجه داشت که مقدمهٔ مقاله دقت تقریباً ۹۰ درصد را در سناریوهای گوناگون و نتیجه‌گیری آن حدود ۹۹ درصد را با آستانه‌های خاص گزارش می‌کند که ناسازگاری درونی محسوب می‌شود \cite{feraudoDIVADIDbasedReputation2024}. مزایای گراف جهت‌دار بی‌دور نسبت به بلاکچین نیز مقیاس‌پذیری، تراکنش‌های سریع و کارمزد پایین گزارش شده است \cite{feraudoDIVADIDbasedReputation2024}.

در طرح \lr{BDRA}\LTRfootnote{\lr{BDRA}}، بلاکچین دولایه‌ای با شناسه‌های غیرمتمرکز به کار رفته است: خودرو شناسهٔ خود را می‌سازد، واحدهای کنارجاده‌ای اعتبارنامه‌های یکتا صادر می‌کنند، تأییدکننده شناسهٔ فرستنده را با فهرست مجاز بررسی کرده و آستانهٔ شهرت را کنترل می‌کند و بازخورد برای به‌روزرسانیِ شهرت به واحد کنارجاده‌ای ارسال می‌شود \cite{mazzocca2025survey}. در طرح دیگری، تغییر آدرس بلاکچین با حفظ شهرت امکان‌پذیر است؛ آدرس‌های جدید از راز مشتق‌شده از کلید خصوصی شناسهٔ غیرمتمرکز ساخته می‌شوند و زنجیره‌ای قطعی می‌سازند که برای حسابرسی یا تحقیق قابل بازسازی است \cite{mazzocca2025survey}. سامانه‌های شهرت معمولاً با سازوکارهای انگیزشی همراه هستند: سامانهٔ پاداش (اعتبار) و سامانهٔ پاداش‌و‌جریمه\LTRfootnote{\lr{Payment-Punishment}} پاداش همکاری و جریمهٔ رفتار خودخواهانه را پیشنهاد می‌کنند \cite{farsimadanReviewSecurityChallenges2025}.

\subsection{اعتماد و حفظ حریم خصوصی}

سامانه‌های اعتماد و شهرت شفافیت رفتار گره‌ها را افزایش می‌دهند، اما این شفافیت می‌تواند با حریم خصوصی رانندگان در تعارض باشد. در تعریف مرور فارسیمدان و همکاران، حریم خصوصی به این معناست که تنها افراد مجاز در شبکهٔ اقتضایی خودرویی بتوانند به داده‌های خودرو، از جمله شناسهٔ واقعی و اطلاعات مکانی، دسترسی و کنترل داشته باشند \cite{farsimadanReviewSecurityChallenges2025}. از سوی دیگر، مشارکت میان خودروها نگرانی‌های حریم خصوصی ایجاد می‌کند \cite{farsimadanReviewSecurityChallenges2025} و کاربران نمی‌خواهند اطلاعات خودرو و مقصدشان به اشتراک گذاشته شود؛ ازاین‌رو تعادل میان منافع عمومی و حریم خصوصی لازم است \cite{hozouriOverviewVANETVehicular2023}. به همین دلیل، پروتکل‌های اعتماد حافظ حریم خصوصی توسعه یافته‌اند \cite{farsimadanReviewSecurityChallenges2025}: در میان طرح‌های مرورشده می‌توان به \lr{PPTM}\LTRfootnote{\lr{PPTM}}، \lr{BTMPP}\LTRfootnote{\lr{BTMPP}} مبتنی بر برش مجموعهٔ خصوصی با فیلتر بلوم، و \lr{PPRM}\LTRfootnote{\lr{PPRM}} و \lr{PPRU}\LTRfootnote{\lr{PPRU}} بر پایهٔ رمزنگاری \lr{ECC} و \lr{Paillier}\LTRfootnote{\lr{Paillier}} با فرض ارائه‌دهندهٔ خدمات ابری «صادق اما کنجکاو» اشاره کرد \cite{farsimadanReviewSecurityChallenges2025}.

در رویکردهای مبتنی بر شناسهٔ غیرمتمرکز، یک خودرو چندین شناسهٔ غیرمتمرکز ثبت می‌کند که همه از طریق اعتبارنامه‌های مبتنی بر اثبات دانش‌صفر\LTRfootnote{\lr{Zero-Knowledge Proof (ZKP)}} صادرشده توسط خودروساز به شناسهٔ اصلی متصل هستند، یا از محاسبهٔ چندطرفه\LTRfootnote{\lr{Multi-Party Computation (MPC)}} استفاده می‌کنند؛ خودروها همچنین می‌توانند با ساختن دو ارائهٔ تأییدپذیر از شبه‌نام\LTRfootnote{\lr{pseudonym}} استفاده کنند \cite{mazzocca2025survey}. در ارتباطات خودرو به خودرو، شناسه‌های غیرمتمرکز و اعتبارنامه‌های تأییدپذیر یکپارچگی، تصدیق هویت، محرمانگی و حریم خصوصی را بدون اتکا به مرجع متمرکز قابل اعتماد تضمین می‌کنند \cite{mazzocca2025survey}.

در طرح‌های مرورشدهٔ دیگر، گروه‌شناسه به جای هویت واقعی به کار رفته است \cite{farsimadanReviewSecurityChallenges2025} و شبه‌نام همراه با امضای آستانه‌ای برای بازیابی هویت خودروی مخرب و تصدیق آستانه‌ای پیشنهاد شده است \cite{farsimadanReviewSecurityChallenges2025}؛ همچنین استفاده از نام مستعار (نشانی عمومی) در بلاکچین با امضای گروهی مبتنی بر هویت، و ردیابی هویت مخرب توسط مرجع مورد اعتماد گزارش شده است \cite{razaBlockchainBasedReputationTrust2024}. استفاده از \lr{zk-SNARKs}\LTRfootnote{\lr{zk-SNARKs}} و اثبات‌های دانش‌صفر غیرتعاملی\LTRfootnote{\lr{Non-Interactive Zero-Knowledge Proof (NIZK)}} نیز برای ناشناس‌ماندن و حریم خصوصی شرطی\LTRfootnote{\lr{conditional privacy}} در شبکه‌های خودرویی گزارش شده است \cite{razaBlockchainBasedReputationTrust2024}. با این وجود، بلاکچین محرمانگی را در میان اهداف اصلی خود ندارد و این یکی از چالش‌های حریم خصوصی در چنین سامانه‌هایی محسوب می‌شود \cite{razaBlockchainBasedReputationTrust2024}.
